\documentclass[10pt,UTF8]{article}


\usepackage{geometry}

\usepackage[scheme=plain]{ctex}

\geometry{a4paper,scale=0.9}

% \usepackage{times}  % 设置正文字体为 Adobe Times Roman
% \usepackage{mathptmx} % 只加载数学字体
% \renewcommand{\rmdefault}{cmr} % 恢复正文字体为 Computer Modern Roman

\usepackage{bm} % 数学粗体（避免斜体被取消）

% \usepackage{makecell}
% \usepackage{multirow} % 表格多行合并

\usepackage{booktabs} % 表格美化

% \usepackage{graphicx} %插入图片的宏包
% \usepackage{subfigure}
% \usepackage{float} %设置图片浮动位置的宏包

\usepackage{amsmath}
\usepackage{esint} % 提供更丰富的积分符号
\usepackage{amssymb}
\usepackage{mathtools}

\usepackage{physics}

% \usepackage{siunitx}
% % 关闭警告
% \ExplSyntaxOn
% \msg_redirect_name:nnn { siunitx } { physics-pkg } { none }
% \ExplSyntaxOff

% \usepackage{bussproofs} % 自然演绎推理树

% 交叉引用
% 关闭边框，设置颜色
\usepackage[colorlinks=true, linkcolor=darkgray, citecolor=teal,
urlcolor=blue]{hyperref}
% 使 \cref 和 \Cref 可用
\usepackage{cleveref}

\usepackage{bookmark}  % 生成 PDF 大纲

\allowdisplaybreaks[3]  % 允许换页


%%%%%%%%%%%%%%%%%%%%%%%% tcolorbox %%%%%%%%%%%%%%%%%%%%%%%%%%
\usepackage[most]{tcolorbox}

\newenvironment{tipBox}
{
  \begin{tcolorbox}[colback=green!5!white,colframe=green!75!black,parbox=false,before
  upper=\par,before lower=\par]}
  {
\end{tcolorbox}}

\newenvironment{conceptBox}
{
\begin{tcolorbox}[colback=blue!5!white,colframe=blue!75!black]}
  {
\end{tcolorbox}}

\newenvironment{theoremBoxNoIndent}
{
\begin{tcolorbox}[colback=yellow!5!white,colframe=yellow!75!black]}
  {
\end{tcolorbox}}

\newenvironment{definitionBox}
{
  \begin{tcolorbox}[colback=blue!5!white,colframe=blue!75!black,parbox=false,before
  upper=\par,before lower=\par]}
  {
\end{tcolorbox}}

\newenvironment{theoremBox}
{
  \begin{tcolorbox}[colback=yellow!5!white,colframe=yellow!75!black,parbox=false,before
  upper=\par,before lower=\par]}
  {
\end{tcolorbox}}
%%%%%%%%%%%%%%%%%%%%%%%%%%%%%%%%%%%%%%%%%%%%%%%%%%%%%%%%%%%%

%%%%%%%%%%%%%%%%%%%%% amsthm %%%%%%%%%%%%%%%%
\usepackage{amsthm}

\theoremstyle{definition}
\newtheorem{definition inner}{Definition}[section]
\newenvironment{definition}[1][]
{
  \begin{definitionBox}
  \begin{definition inner}[#1]\hfill\par}
    {
    \end{definition inner}
\end{definitionBox}}

\theoremstyle{definition}
\newtheorem{theorem inner}{Theorem}[section]
\newenvironment{theorem}[1][]
{
  \begin{theoremBox}
  \begin{theorem inner}[#1]\hfill\par}
    {
    \end{theorem inner}
\end{theoremBox}}

\theoremstyle{definition}
\newtheorem*{proof inner}{\textit{Proof}}
\renewenvironment{proof}[1][]
{
\begin{proof inner}[#1]\hfill\par}
  {\qed\vphantom{.}
\end{proof inner}}

\theoremstyle{plain}
\newtheorem{tip inner}{Tip}[section]
\newenvironment{tip}[1][]
{
  \begin{tipBox}
  \begin{tip inner}[#1]\hfill\par}
    {
    \end{tip inner}
\end{tipBox}}

\theoremstyle{remark}
\newtheorem*{remark}{\textit{Remark}}
%%%%%%%%%%%%%%%%%%%%%%%%%%%%%%%%%%%%%%%%%%%%%%


\newcommand{\iu}{\mathrm{i}}  % 虚数单位 i


\begin{document}

\part*{Higher Order ODEs}

\section{Linear ODEs}

\subsection{General Theory}

\begin{theorem}[Existence and Uniqueness Theorem]
  {}\label{thm: Existence and Uniqueness}
  Let \(I \subset\mathbb{R}\) be an open interval and suppose
  functions \(g, p_0, p_1, \ldots, p_{n-1}\) are continuous on \(I\).
  Then for \(t_0 \in I\) and real numbers \(y_0, y_1, \ldots, y_{n-1}\),
  there exists exactly one solution to the IVP
  \begin{equation*}
    \begin{cases}
      y^{(n)} + p_{n-1}(t) y^{(n-1)} + \cdots + p_1(t) y' + p_0(t) y = g(t), \\
      y(t_0) = y_0, \\
      y'(t_0) = y_1, \\
      \vdots \\
      y^{(n-1)}(t_0) = y_{n-1},
    \end{cases}
  \end{equation*}
\end{theorem}

\begin{theorem}[Principle of Superposition]
  {}\label{thm: Principle of Superposition}
  If \(y_1, y_2, \ldots, y_n\) are solutions to the homogeneous linear ODE
  \begin{equation*}
    y^{(n)} + p_{n-1}(t) y^{(n-1)} + \cdots + p_1(t) y' + p_0(t) y = 0,
  \end{equation*}
  then any linear combination
  \begin{equation*}
    y = c_1 y_1 + c_2 y_2 + \cdots + c_n y_n,
  \end{equation*}
  where \(c_1, c_2, \ldots, c_n\) are constants, is also a solution
  to the homogeneous equation.
\end{theorem}

\begin{definition}[Wronskian]
  {}\label{def: Wronskian}
  The Wronskian of \(n\) functions \(y_1, y_2, \ldots, y_n\) is
  defined as the determinant
  \begin{equation*}
    W(y_1, y_2, \ldots, y_n)(t) =
    \begin{vmatrix}
      y_1(t) & y_2(t) & \cdots & y_n(t) \\
      y_1'(t) & y_2'(t) & \cdots & y_n'(t) \\
      \vdots & \vdots & \ddots & \vdots \\
      y_1^{(n-1)}(t) & y_2^{(n-1)}(t) & \cdots & y_n^{(n-1)}(t)
    \end{vmatrix}.
  \end{equation*}
\end{definition}

\begin{definition}[Fundamental Set of Solutions]
  {}\label{def: Fundamental Set of Solutions}
  The set of \(n\) solutions \(y_1, y_2, \ldots, y_n\) to
  an \(n\)-th order homogeneous linear ODE defined on an interval \(I\)
  is called a \textbf{fundamental set of solutions (FSS)}
  if every solution to the homogeneous equation can be
  expressed as a linear combination of \(y_1, y_2, \ldots, y_n\).
\end{definition}

\begin{theorem}
  If \(p_0,\ldots,p_{n-1}\) are continuous on an interval \(I\) and
  \(y_1, \ldots, y_n\) are solutions to the homogeneous linear ODE
  \begin{equation*}
    y^{(n)} + p_{n-1}(t) y^{(n-1)} + \cdots + p_1(t) y' + p_0(t) y = 0,
  \end{equation*}
  then every solution to the ODE can be expressed as a linear
  combination of \(y_1, \ldots, y_n\) if and only if
  \begin{equation*}
    \exists t\in I \text{ s.t. } W(y_1, \ldots, y_n)(t) \neq 0
  \end{equation*}
\end{theorem}

\begin{theorem}
  Let \(y_1,\ldots,y_n\) be the solutions to the homogeneous linear ODE
  \begin{equation*}
    y^{(n)} + p_{n-1}(t) y^{(n-1)} + \cdots + p_1(t) y' + p_0(t) y = 0,
    \qfor t\in I.
  \end{equation*}
  Then
  \begin{equation*}
    y_1, \ldots, y_n \text{ are linearly independent}
    \quad\iff\quad
    \exists t\in I \text{ s.t. } W(y_1, \ldots, y_n)(t) \neq 0.
  \end{equation*}
\end{theorem}

\begin{theorem}[Abel's Theorem]
  {}\label{thm: Abel's Theorem}
  If \(y_1, \ldots, y_n\) are solutions to the homogeneous linear ODE
  \begin{equation*}
    y^{(n)} + \textcolor{blue}{p_{n-1}(t)} y^{(n-1)}
    + p_{n-2}(t) y^{(n-2)} + \cdots + p_1(t) y' + p_0(t) y = 0,
    \qfor t\in I,
  \end{equation*}
  then the Wronskian is given by
  \begin{equation*}
    W(y_1, \ldots, y_n)(t) = \textcolor{red}{c} \cdot
    \exp(-\int \textcolor{blue}{p_{n-1}(t)} \dd{t})
    \qquad\text{ where } \textcolor{red}{c}
    \text{ is a constant depending on } y_1, \ldots, y_n.
  \end{equation*}
\end{theorem}

\subsection{Homogeneous Equations with Constant Coefficients}

For constants \(a_n \ne 0, a_{n-1}, \ldots, a_1, a_0 \in \mathbb{R}\),
\begin{equation*}
  a_n y^{(n)} + a_{n-1} y^{(n-1)} + \cdots + a_1 y' + a_0 y = 0
\end{equation*}
\textbf{Characteristic polynomial}:
\begin{align*}
  Z(r) &= a_n r^n + a_{n-1} r^{n-1} + \cdots + a_1 r + a_0 \\
  &= a_n (r - r_1)(r - r_2) \cdots (r - r_n) \\
  &= a_n {(r - r_1)}^{m_1} {(r - r_2)}^{m_2} \cdots {(r - r_k)}^{m_k}
\end{align*}
where \(r_1, r_2, \ldots, r_n \in \mathbb{C}\) are the roots of \(Z(r)\).

\subsubsection*{Complex-valued solutions}

Complex roots must appear in conjugate pairs, i.e., if
\(r = \lambda + \iu \mu \) is a root,
then its conjugate \(\bar{r} = \lambda - \iu \mu \) is also a root.
The complex-valued solution pair \(y = e^{(\lambda + \iu \mu) t}\)
and \(y = e^{(\lambda - \iu \mu) t}\) can be re-expressed
as the real-valued solution pair:
\begin{equation*}
  y = e^{\lambda t} \cos(\mu t) \qand
  y = e^{\lambda t} \sin(\mu t).
\end{equation*}

\subsection{Non-Homogeneous Equations}



\subsection{Method of Undetermined Coefficients}


\subsection{Variation of Parameters}

\end{document}
